\documentclass[11pt]{article}
\usepackage{mathunal-base}
\mathunalfuente{serif}

\renewcommand{\examtitulo}{Métodos Numéricos (3006907) --- Supletorio, Segundo Examen Parcial}
\renewcommand{\examfecha}{Semestre 02-2016, 31 de Octubre de 2016}

\begin{document}

\examheader

\vspace{4pt}
\noindent\textit{Duración: 1 hora y 50 minutos.}

\vspace{6pt}
\noindent Nombre completo \dotfill

\noindent Carné \dotfill\ Grupo \dotfill

\vspace{8pt}
\noindent\textbf{La interpretación del examen hace parte de la evaluación, por tanto NO se admiten preguntas una vez comenzada la prueba.}

\vspace{10pt}
\noindent\textbf{1. (15\%)} Encuentre el ajuste para una curva de la forma $y=\dfrac{B}{A+x}$ para los siguientes conjuntos de datos, mediante la teoría de los mínimos cuadrados.
\[
\begin{array}{c|c|c|c|c|c}
x_i & -1 & 0 & 1 & 2 & 5\\ \hline
y_i & -\dfrac12 & \dfrac14 & \dfrac15 & \dfrac27 & \dfrac53
\end{array}
\]

\vspace{5mm}
\noindent\textbf{2. (15\%)} Considere la siguiente tabla de diferencias divididas
\[
\begin{array}{c|c|c|c}
x_k & f[x_k] & f[x_{k-1},x_k] & f[x_{k-2},x_{k-1},x_k]\\ \hline
x_0=-1 & \gamma & & \\
x_1=1 & \alpha & -3 & \\
x_2=2.5 & 4 & \beta & 4
\end{array}
\]
Encuentre el valor de $\alpha=f(1)$.

\vspace{5mm}
\noindent\textbf{3. (25\%)} Considere la función $S(x)$ definida por:
\[
S(x)=\begin{cases}
1+a(x-1)+b(x-1)^3, & 1\leq x\leq 2\\
1+c(x-2)-\dfrac34(x-2)^2+d(x-2)^3, & 2\leq x\leq 3
\end{cases}
\]
\begin{enumerate}
\item[a.] (20\%) Determine los valores de $a,b,c$ y $d$ que hacen que $S(x)$ sea un spline cúbico natural para el conjunto de puntos $(x,f(x))$ dado por $\{(1,1),(2,1),(3,0)\}$.
\item[b.] (5\%) Obtenga el valor aproximado de $f(3/2)$.
\end{enumerate}

\vspace{5mm}
\noindent\textbf{4. (25\%)} Considere la fórmula de cuadratura
\[
\int_0^1 f(x)\,dx\approx Af(0)+Bf(z)
\]
\begin{enumerate}
\item[i.] (15\%) Calcule los coeficientes $A$ y $B$ y el nodo $z$ de tal manera que la fórmula sea exacta para todos los polinomios de grado tan alto como sea posible. ¿Cuál es el grado de precisión de esta fórmula?
\item[ii.] (10\%) Use la fórmula anterior para dar un valor aproximado de $\displaystyle\int_{\pi/4}^{\pi/3}\cos x^2\,dx$.
\end{enumerate}

\vspace{5mm}
\noindent\textbf{5. (20\%)} Las \textbf{fórmulas de Newton-Cotes abiertas} no incluyen los extremos del intervalo $[a,b]$. Éstas utilizan los nodos $x_j=a+(j+1)h$, para $j=-1,0,\ldots,n,n+1$ donde $h=\dfrac{b-a}{n+2}$. Esto implica que los nodos interiores serán los $x_i$ para $i=0,\ldots,n$ y los extremos del intervalo son $x_{-1}=a$ y $x_{n+1}=b$.

La fórmula de Newton-Cotes abierta que se obtiene al tomar 2 subintervalos, es decir, al tomar $h=\dfrac{b-a}2$ se conoce como la \textit{regla del punto medio simple} y esta dada por:
\[
\int_a^b f(x)\,dx=\int_{x_{-1}}^{x_1} f(x)\,dx\approx 2hf(x_0).
\]
\textbf{Deduzca} la \underline{regla del punto medio compuesta} que se obtiene al tomar $h=\dfrac{b-a}{n+2}$, nodos $x_j=a+(j+1)h$, $j=-1,0,\ldots,n,n+1$, para $n$ entero par mayor o igual a 2.

\vfill
\rule{\linewidth}{0.4pt}\\[1pt]
{\scriptsize\textit{Transcripción digital --- MathUNAL}\hfill\textcolor{unalblue}{\textbf{1}}}

\end{document}
